\begin{pastexam}[label={px:ntut-114-ee-ds-11}]{北科大電機所（戊組） 114 學年度 第 III-3 題}
（追蹤＋問答；20 分）

You have a hash table with 19 slots and the following hash function: $h(k)=(k+10)\ mod\ 19$. Insert
the following keys in order: 10, 29, 47, 38, 24, 5, 28, 9, 17.
\begin{enumerate}[label=(\alph*)]
\item How many collisions occurred when using chaining versus linear probing for collision resolution?
Compare and explain the difference. \textbf{(6\%)}
\item If you want to search for key 19, analyze its possible position(s) under both chaining and linear
probing methods. Which method would be more efficient for this search and why? \textbf{(6\%)}
\item The density of the hash table is 47.4\%. Explain how this density affects the search performance
in chaining and linear probing methods. Which method would you recommend for this density and why?
\textbf{(8\%)}
\end{enumerate}
\end{pastexam}

\begin{pastsoln}
先算每個鍵的雜湊值 $h(k) = (k + 10) \bmod 19$：
\begin{center}
\small
\begin{tabular}{l*{9}{c}}
\toprule
$k$ & 10 & 29 & 47 & 38 & 24 & 5 & 28 & 9 & 17 \\
$k + 10$ & 20 & 39 & 57 & 48 & 34 & 15 & 38 & 19 & 27 \\
$h(k)$ & 1 & 1 & 0 & 10 & 15 & 15 & 0 & 0 & 8 \\
\bottomrule
\end{tabular}
\end{center}

\textbf{(a)} 「碰撞」算成：要放的位置已經有別的鍵。

\textbf{鏈結法}（資結書第~12.2~節）：每格是一條串列，鍵一定放在自己的那格。 放進時那格已經有鍵的是
$29$（第 1 格有 10）、$5$（第 15 格有 24）、$28$（第 0 格有 47）、$9$（第 0 格有 47、28），共 \textbf{4 次}。
最後：第 0 格 $47, 28, 9$；第 1 格 $10, 29$；第 8 格 $17$；第 10 格 $38$；第 15 格 $24, 5$。

\textbf{線性探測}（資結書第~12.3~節）：位置被占就往下一格找（第 18 格的下一格是第 0 格）。
\begin{center}
\small
\begin{tabular}{rcll}
\toprule
$k$ & $h(k)$ & 探測經過的格子 & 碰撞次數 \\
\midrule
10 & 1 & 1 & 0 \\
29 & 1 & 1（有 10）$\to$ 2 & 1 \\
47 & 0 & 0 & 0 \\
38 & 10 & 10 & 0 \\
24 & 15 & 15 & 0 \\
5 & 15 & 15（有 24）$\to$ 16 & 1 \\
28 & 0 & 0（47）$\to$ 1（10）$\to$ 2（29）$\to$ 3 & 3 \\
9 & 0 & 0 $\to$ 1 $\to$ 2 $\to$ 3（28）$\to$ 4 & 4 \\
17 & 8 & 8 & 0 \\
\bottomrule
\end{tabular}
\end{center}
探測到已占用的格子共 $1 + 1 + 3 + 4 =$ \textbf{9 次}（發生過碰撞的鍵是 4 個，和鏈結法一樣）。 最後的表：
第 0--4 格依序是 $47, 10, 29, 28, 9$，第 8 格 $17$，第 10 格 $38$，第 15、16 格 $24, 5$，其餘空著。

\textbf{差別在哪。} 鏈結法只會和「雜湊值相同」的鍵碰撞。 線性探測的鍵會占走別人的家：$29$ 放到第 2 格，
$28$ 的起點雖然是第 0 格，卻得一路經過 $10$、$29$（它們的雜湊值是 1）；第 0--4 格連成一大段，任何落在
這一段的鍵都要走到段尾，段又再變長。 這就是主要群聚（primary clustering），所以同樣的資料，線性
探測的碰撞次數比較多（9 比 4）。

\textbf{(b)} $h(19) = 29 \bmod 19 = 10$。
\begin{itemize}
\item 鏈結法：$19$ 如果在表裡，只可能在第 10 格的串列裡。 那條串列只有 $38$，比較 $1$ 次就知道 $19$ 不在表裡。
\item 線性探測：$19$ 可能在第 10 格，或是從第 10 格往後、第一個空格之前的任何一格。 第 10 格是 $38$，不是；
第 11 格是空的，搜尋停止，$19$ 不在表裡。 共探測 $2$ 格（比較一個鍵、看到一個空格）。
\end{itemize}
這次鏈結法比較有效率（比較 1 個鍵，線性探測要看 2 格）。 原因是鏈結法只看雜湊值相同的鍵；線性探測要
一直走到空格才能確定找不到，若起點落在第 0--4 格那種長段裡，要走更多格（例如找一個雜湊值是 $0$ 的
不存在的鍵，要看第 0--5 格共 6 格）。

\textbf{(c)} 密度（負載因子）$\alpha = 9/19 \approx 0.474$。
\begin{itemize}
\item 鏈結法：串列的平均長度是 $\alpha$。 失敗搜尋平均比較 $\alpha \approx 0.47$ 個鍵；成功搜尋平均比較
$1 + \frac{n - 1}{2m} = 1 + \frac{8}{38} \approx 1.21$ 個鍵（資結書定理~12.2）。 效能隨 $\alpha$ 線性變差，
$\alpha$ 超過 $1$ 也還能用。 本題這張表實際的成功搜尋平均是 $(1 + 2 + 3 + 1 + 2 + 1 + 1 + 1 + 2)/9 = 14/9 \approx 1.56$。
\item 線性探測：照 Knuth 的近似（資結書第~12.4~節），成功搜尋約
$\frac{1}{2}\bigl(1 + \frac{1}{1 - \alpha}\bigr) = \frac{1}{2}(1 + 1.9) = 1.45$ 次探測，失敗搜尋約
$\frac{1}{2}\bigl(1 + \frac{1}{(1 - \alpha)^2}\bigr) = \frac{1}{2}(1 + 3.61) \approx 2.31$ 次。 $\alpha$ 越接近 $1$，
$\frac{1}{(1-\alpha)^2}$ 爆增，而且表滿了就放不下。 本題實際的成功搜尋平均是
$(1 + 2 + 1 + 1 + 1 + 2 + 4 + 5 + 1)/9 = 18/9 = 2$ 次，已經看得到第 0--4 格的群聚。
\end{itemize}
\textbf{建議用鏈結法。} 在 $\alpha \approx 0.47$ 時兩種方法都還算快，但這組資料雜湊值集中在 $0$、$1$，線性
探測已經形成 5 格長的群聚，搜尋成本比鏈結法高（成功搜尋平均 2 比 1.56，碰撞 9 比 4），再插入更多鍵會
惡化得更快。 鏈結法只受同一格的鍵影響，效能穩定，刪除也簡單（線性探測刪除要用墓碑）。 線性探測的
優點是不需要指標、記憶體連續、快取友善；若記憶體很緊、而且能保證 $\alpha$ 維持在 $0.5$ 以下、雜湊值分布
均勻，也可以選它。
\end{pastsoln}
